\scanpage{42}\mspage{48}
On applique Dér 2, Dér 4 (g), pour se ramener à regarder, \inserted{pour $y_0\in Y$}, l'hom $\alpha$ du diagramme \inserted{(commutatif !)} suivant :
\diagram{\begin{tikzcd}[ampersand replacement=\&,column sep=large]
 \Gamma(\{y_0\},F_{y_0})\unsure{=|\leftarrow|\simeq}F_{y_0}\arrow[r,"\alpha"] \& \Gamma_{\mathbb D}(X_{/y_0},f^*(F))\\
 \Gamma(Y_{/y_0},F)\arrow[u,leftrightarrow,"\varphi","\unsure{\sim|?}"']\arrow[ur,"\alpha'"'] \&
\end{tikzcd}}
\margin{$\longleftarrow$ \deleted{défini par le foncteur $\unsure{\{y\}|f_y}$}}
où $\alpha'$ est déduit du foncteur induit par $f$
\[
 X_{/y_0}\longrightarrow Y_{/y_0}
\]
et où $\varphi$ est déduit du foncteur
\[
 \{y_0\}\xrightarrow[\substack{(\text{objet}\\ \text{final})}]{} Y_{/y_0}
\]

% slanted
\margin{Cf cor. 3 page 50, qui devrait remonter ici.}
Donc $\varphi$ est un iso par Der 5 (g)$^*$ (NB Si $z\in Z$ d'où $\{z\}\to Z$, ce foncteur est coh. cofinal ssi $z$ est objet final, \deleted{\unsure{donc $f$}} \unsure{hom.|homo} final ssi $z$ est objet initial) Donc la conclusion du corollaire équivaut à dire que $\alpha'$ est iso. Mais $X_{/y_0}\to Y_{/y_0}$ est aussi hom. cofinal \deleted{\unsure{\(\longrightarrow\)}} ; car si $y_1\xrightarrow{u}y_0$ \deleted{\unsure{objet}} $\in$ de $Y_{/y_0}$, on a un iso. can.
\[
 (X_{/y_0})_{/(y_1,u)}\simeq X_{/y_1}\qquad\text{qui est contractile}
\]
Donc par la \uline{deuxième partie} de \deleted{Dér} (5)(g), on gagne.
% slanted
\margin{qui est \uline{fausse}}
% archive sheet number: 41
